% Lines 3813-4399 from original attention_transformers_notes_ghost.tex
% Chapter 13: Domain-Specific Applications

\section{Domain-Specific Applications}
%==============================================================================

\subsection{Transformers for Computer Vision}

\subsubsection{Vision Transformer (ViT) Implementation}

Complete implementation of Vision Transformer:

\begin{lstlisting}[language=Python, caption=Vision Transformer Implementation]
@dataclass(frozen=True, slots=True)
class ViTConfig:
    image_size: int = 224
    patch_size: int = 16
    num_channels: int = 3
    num_classes: int = 1000
    d_model: int = 768
    n_heads: int = 12
    n_layers: int = 12
    d_ff: int = 3072
    dropout: float = 0.1
    attention_dropout: float = 0.0

class PatchEmbedding(nn.Module):
    """Convert image into sequence of patch embeddings"""

    def __init__(self, config: ViTConfig) -> None:
        super().__init__()
        self.config = config
        self.num_patches = (config.image_size // config.patch_size) ** 2

        # Linear projection of flattened patches
        self.projection = nn.Linear(
            config.num_channels * config.patch_size ** 2,
            config.d_model
        )

    def forward(self, x: Tensor) -> Tensor:
        # x shape: (batch, channels, height, width)
        batch_size = x.shape[0]

        # Create patches
        # (batch, channels, h, w) -> (batch, num_patches, patch_dim)
        patches = x.unfold(2, self.config.patch_size, self.config.patch_size) \
                   .unfold(3, self.config.patch_size, self.config.patch_size) \
                   .contiguous()

        patches = patches.view(
            batch_size,
            self.config.num_channels,
            self.num_patches,
            self.config.patch_size ** 2
        )
        patches = patches.permute(0, 2, 1, 3).contiguous()
        patches = patches.view(
            batch_size,
            self.num_patches,
            -1
        )

        # Project patches
        embeddings = self.projection(patches)
        return embeddings

class VisionTransformer(nn.Module):
    """Vision Transformer for image classification"""

    def __init__(self, config: ViTConfig) -> None:
        super().__init__()
        self.config = config

        # Patch embedding
        self.patch_embed = PatchEmbedding(config)
        num_patches = self.patch_embed.num_patches

        # Class token
        self.cls_token = nn.Parameter(
            torch.zeros(1, 1, config.d_model)
        )

        # Position embeddings
        self.pos_embed = nn.Parameter(
            torch.zeros(1, num_patches + 1, config.d_model)
        )

        # Dropout
        self.pos_drop = nn.Dropout(config.dropout)

        # Transformer encoder
        encoder_layer = nn.TransformerEncoderLayer(
            d_model=config.d_model,
            nhead=config.n_heads,
            dim_feedforward=config.d_ff,
            dropout=config.dropout,
            activation='gelu',
            batch_first=True
        )
        self.encoder = nn.TransformerEncoder(
            encoder_layer,
            num_layers=config.n_layers
        )

        # Classification head
        self.norm = nn.LayerNorm(config.d_model)
        self.head = nn.Linear(config.d_model, config.num_classes)

        # Initialize weights
        self._init_weights()

    def _init_weights(self) -> None:
        # Initialize patch embeddings
        nn.init.xavier_uniform_(self.patch_embed.projection.weight)
        nn.init.normal_(self.patch_embed.projection.bias, std=1e-6)

        # Initialize positional embeddings
        nn.init.normal_(self.pos_embed, std=0.02)
        nn.init.normal_(self.cls_token, std=0.02)

        # Initialize classification head
        nn.init.zeros_(self.head.bias)
        nn.init.xavier_uniform_(self.head.weight)

    def forward(self, x: Tensor) -> Tensor:
        batch_size = x.shape[0]

        # Patch embedding
        x = self.patch_embed(x)

        # Add class token
        cls_tokens = self.cls_token.expand(batch_size, -1, -1)
        x = torch.cat([cls_tokens, x], dim=1)

        # Add position embeddings
        x = x + self.pos_embed
        x = self.pos_drop(x)

        # Transformer encoder
        x = self.encoder(x)

        # Classification head uses only class token
        x = self.norm(x[:, 0])
        x = self.head(x)

        return x

    def get_attention_maps(
        self,
        x: Tensor
    ) -> list[Tensor]:
        """Extract attention maps for visualization"""

        batch_size = x.shape[0]

        # Patch embedding
        x = self.patch_embed(x)

        # Add class token
        cls_tokens = self.cls_token.expand(batch_size, -1, -1)
        x = torch.cat([cls_tokens, x], dim=1)

        # Add position embeddings
        x = x + self.pos_embed

        attention_maps = []

        # Extract attention from each layer
        for layer in self.encoder.layers:
            # Get attention weights
            with torch.no_grad():
                _, attn_weights = layer.self_attn(
                    x, x, x,
                    need_weights=True,
                    average_attn_weights=False
                )
            attention_maps.append(attn_weights)

            # Forward through layer
            x = layer(x)

        return attention_maps
\end{lstlisting}

\subsubsection{Hybrid CNN-Transformer Architectures}

Combine convolutional layers with transformers:

\begin{lstlisting}[language=Python, caption=Hybrid CNN-Transformer]
class ConvStem(nn.Module):
    """Convolutional stem for hybrid architectures"""

    def __init__(
        self,
        in_channels: int,
        out_channels: int,
        patch_size: int
    ) -> None:
        super().__init__()

        # Progressive downsampling with conv layers
        self.layers = nn.Sequential(
            # Initial conv
            nn.Conv2d(in_channels, out_channels // 4, 3, 2, 1),
            nn.BatchNorm2d(out_channels // 4),
            nn.ReLU(inplace=True),

            # Second conv
            nn.Conv2d(out_channels // 4, out_channels // 2, 3, 2, 1),
            nn.BatchNorm2d(out_channels // 2),
            nn.ReLU(inplace=True),

            # Final conv
            nn.Conv2d(out_channels // 2, out_channels, 3, 2, 1),
            nn.BatchNorm2d(out_channels),
            nn.ReLU(inplace=True),

            # Adjust to patch size
            nn.Conv2d(out_channels, out_channels, patch_size // 8, patch_size // 8),
        )

    def forward(self, x: Tensor) -> Tensor:
        x = self.layers(x)
        # Flatten spatial dimensions
        batch_size, channels, h, w = x.shape
        x = x.view(batch_size, channels, h * w).transpose(1, 2)
        return x

class HybridVisionTransformer(nn.Module):
    """Vision Transformer with CNN stem"""

    def __init__(
        self,
        config: ViTConfig,
        use_conv_stem: bool = True
    ) -> None:
        super().__init__()
        self.config = config

        if use_conv_stem:
            self.patch_embed = ConvStem(
                config.num_channels,
                config.d_model,
                config.patch_size
            )
            # Calculate number of patches after conv stem
            self.num_patches = (config.image_size // config.patch_size) ** 2
        else:
            self.patch_embed = PatchEmbedding(config)
            self.num_patches = self.patch_embed.num_patches

        # Rest is similar to standard ViT
        self.cls_token = nn.Parameter(
            torch.zeros(1, 1, config.d_model)
        )
        self.pos_embed = nn.Parameter(
            torch.zeros(1, self.num_patches + 1, config.d_model)
        )

        # Transformer encoder
        self.encoder = nn.TransformerEncoder(
            nn.TransformerEncoderLayer(
                d_model=config.d_model,
                nhead=config.n_heads,
                dim_feedforward=config.d_ff,
                dropout=config.dropout,
                activation='gelu',
                batch_first=True
            ),
            num_layers=config.n_layers
        )

        self.norm = nn.LayerNorm(config.d_model)
        self.head = nn.Linear(config.d_model, config.num_classes)

    def forward(self, x: Tensor) -> Tensor:
        batch_size = x.shape[0]

        # Patch embedding (CNN or linear)
        x = self.patch_embed(x)

        # Add class token
        cls_tokens = self.cls_token.expand(batch_size, -1, -1)
        x = torch.cat([cls_tokens, x], dim=1)

        # Add position embeddings
        x = x + self.pos_embed

        # Transformer encoder
        x = self.encoder(x)

        # Classification head
        x = self.norm(x[:, 0])
        x = self.head(x)

        return x
\end{lstlisting}

\subsection{Transformers for Time Series}

\subsubsection{Time Series Transformer}

Adapt transformers for time series forecasting:

\begin{lstlisting}[language=Python, caption=Time Series Transformer]
@dataclass(frozen=True, slots=True)
class TimeSeriesConfig:
    seq_len: int  # Input sequence length
    pred_len: int  # Prediction length
    d_input: int  # Number of input features
    d_model: int = 512
    n_heads: int = 8
    n_layers: int = 6
    d_ff: int = 2048
    dropout: float = 0.1

class TimeEmbedding(nn.Module):
    """Embed time features for time series"""

    def __init__(self, d_model: int) -> None:
        super().__init__()

        # Embeddings for different time scales
        self.minute_embed = nn.Embedding(60, d_model)
        self.hour_embed = nn.Embedding(24, d_model)
        self.day_embed = nn.Embedding(32, d_model)
        self.month_embed = nn.Embedding(13, d_model)
        self.weekday_embed = nn.Embedding(7, d_model)

    def forward(self, timestamps: Tensor) -> Tensor:
        # timestamps shape: (batch, seq_len, 5)
        # Contains: [minute, hour, day, month, weekday]

        minute_emb = self.minute_embed(timestamps[..., 0].long())
        hour_emb = self.hour_embed(timestamps[..., 1].long())
        day_emb = self.day_embed(timestamps[..., 2].long())
        month_emb = self.month_embed(timestamps[..., 3].long())
        weekday_emb = self.weekday_embed(timestamps[..., 4].long())

        # Combine embeddings
        time_emb = minute_emb + hour_emb + day_emb + month_emb + weekday_emb
        return time_emb

class TimeSeriesTransformer(nn.Module):
    """Transformer for time series forecasting"""

    def __init__(self, config: TimeSeriesConfig) -> None:
        super().__init__()
        self.config = config

        # Input projection
        self.input_projection = nn.Linear(config.d_input, config.d_model)

        # Time embeddings
        self.time_embed = TimeEmbedding(config.d_model)

        # Positional encoding
        self.pos_encoding = nn.Parameter(
            torch.zeros(1, config.seq_len + config.pred_len, config.d_model)
        )

        # Transformer encoder for input sequence
        encoder_layer = nn.TransformerEncoderLayer(
            d_model=config.d_model,
            nhead=config.n_heads,
            dim_feedforward=config.d_ff,
            dropout=config.dropout,
            batch_first=True
        )
        self.encoder = nn.TransformerEncoder(
            encoder_layer,
            num_layers=config.n_layers
        )

        # Transformer decoder for predictions
        decoder_layer = nn.TransformerDecoderLayer(
            d_model=config.d_model,
            nhead=config.n_heads,
            dim_feedforward=config.d_ff,
            dropout=config.dropout,
            batch_first=True
        )
        self.decoder = nn.TransformerDecoder(
            decoder_layer,
            num_layers=config.n_layers
        )

        # Output projection
        self.output_projection = nn.Linear(config.d_model, config.d_input)

        # Learnable start token for decoder
        self.start_token = nn.Parameter(
            torch.zeros(1, 1, config.d_model)
        )

    def forward(
        self,
        x: Tensor,  # (batch, seq_len, d_input)
        timestamps: Tensor,  # (batch, seq_len + pred_len, 5)
        future_timestamps: Tensor | None = None
    ) -> Tensor:
        batch_size = x.shape[0]

        # Project input
        x = self.input_projection(x)

        # Add time embeddings for input
        input_time_emb = self.time_embed(timestamps[:, :self.config.seq_len])
        x = x + input_time_emb

        # Add positional encoding
        x = x + self.pos_encoding[:, :self.config.seq_len]

        # Encode input sequence
        memory = self.encoder(x)

        # Prepare decoder input
        if self.training:
            # During training, use teacher forcing
            # Shift target sequence and prepend start token
            decoder_input = torch.cat([
                self.start_token.expand(batch_size, -1, -1),
                memory[:, -self.config.pred_len:, :]
            ], dim=1)
        else:
            # During inference, generate autoregressively
            decoder_input = self.start_token.expand(batch_size, -1, -1)

        # Add time embeddings for future
        future_time_emb = self.time_embed(
            timestamps[:, self.config.seq_len:]
        )

        # Generate future predictions
        if self.training:
            # Parallel decoding during training
            decoder_input = decoder_input + future_time_emb
            decoder_input = decoder_input + self.pos_encoding[
                :, self.config.seq_len:self.config.seq_len + self.config.pred_len
            ]

            # Create causal mask
            tgt_mask = nn.Transformer.generate_square_subsequent_mask(
                self.config.pred_len
            ).to(x.device)

            # Decode
            output = self.decoder(
                decoder_input,
                memory,
                tgt_mask=tgt_mask
            )
        else:
            # Autoregressive decoding during inference
            outputs = []
            for i in range(self.config.pred_len):
                # Add time embedding for current step
                pos_emb = self.pos_encoding[
                    :, self.config.seq_len + i:self.config.seq_len + i + 1
                ]
                time_emb = future_time_emb[:, i:i+1]

                decoder_input_step = decoder_input + pos_emb + time_emb

                # Decode one step
                output = self.decoder(decoder_input_step, memory)
                outputs.append(output[:, -1:])

                # Append to decoder input for next step
                if i < self.config.pred_len - 1:
                    decoder_input = torch.cat([
                        decoder_input,
                        output[:, -1:]
                    ], dim=1)

            output = torch.cat(outputs, dim=1)

        # Project to output dimension
        predictions = self.output_projection(output)

        return predictions
\end{lstlisting}

\subsubsection{Temporal Fusion Transformer}

Advanced architecture for interpretable time series forecasting:

\begin{lstlisting}[language=Python, caption=Temporal Fusion Transformer Components]
class GatedResidualNetwork(nn.Module):
    """Gated Residual Network for feature processing"""

    def __init__(
        self,
        input_dim: int,
        hidden_dim: int,
        output_dim: int,
        dropout: float = 0.1
    ) -> None:
        super().__init__()

        self.input_dim = input_dim
        self.output_dim = output_dim

        # Layers
        self.fc1 = nn.Linear(input_dim, hidden_dim)
        self.fc2 = nn.Linear(hidden_dim, output_dim)

        # Gating mechanism
        self.gate_fc1 = nn.Linear(input_dim, hidden_dim)
        self.gate_fc2 = nn.Linear(hidden_dim, output_dim)

        # Skip connection
        if input_dim != output_dim:
            self.skip = nn.Linear(input_dim, output_dim)
        else:
            self.skip = nn.Identity()

        self.elu = nn.ELU()
        self.dropout = nn.Dropout(dropout)
        self.sigmoid = nn.Sigmoid()

    def forward(self, x: Tensor) -> Tensor:
        # Main path
        hidden = self.elu(self.fc1(x))
        hidden = self.dropout(hidden)
        hidden = self.fc2(hidden)

        # Gating path
        gate = self.elu(self.gate_fc1(x))
        gate = self.dropout(gate)
        gate = self.sigmoid(self.gate_fc2(gate))

        # Apply gating and skip connection
        return gate * hidden + (1 - gate) * self.skip(x)

class VariableSelectionNetwork(nn.Module):
    """Select relevant variables/features"""

    def __init__(
        self,
        input_dim: int,
        num_vars: int,
        hidden_dim: int,
        dropout: float = 0.1
    ) -> None:
        super().__init__()

        # Individual variable transformations
        self.var_transforms = nn.ModuleList([
            GatedResidualNetwork(
                input_dim, hidden_dim, hidden_dim, dropout
            )
            for _ in range(num_vars)
        ])

        # Variable selection weights
        self.var_selection = GatedResidualNetwork(
            hidden_dim * num_vars,
            hidden_dim,
            num_vars,
            dropout
        )

        self.softmax = nn.Softmax(dim=-1)

    def forward(self, vars_list: list[Tensor]) -> tuple[Tensor, Tensor]:
        # Transform each variable
        transformed = []
        for i, var in enumerate(vars_list):
            transformed.append(self.var_transforms[i](var))

        # Concatenate transformed variables
        concatenated = torch.cat(transformed, dim=-1)

        # Compute selection weights
        weights = self.var_selection(concatenated)
        weights = self.softmax(weights)

        # Apply weighted selection
        selected = torch.stack(transformed, dim=-1)
        output = torch.sum(
            selected * weights.unsqueeze(-2),
            dim=-1
        )

        return output, weights
\end{lstlisting}

%==============================================================================
